\documentclass[11pt]{article}
\usepackage{amsmath,amssymb,hyperref}
\title{Where to keep the spare: recovery certificates for a tree of shared suppliers}
\author{}
\date{}
\begin{document}
\maketitle

A supplier loss can require changing every assignment along a long chain, even if it changes only one input to the scheduler. But in the model below, ruling out every place to keep the spare may require just two fault cases. That guarantee depends on which sites the operator permits: if the permitted sites form $c$ connected components, the bound is $\max(2,c)$ and $c$ can be necessary. Finding the smallest explanation in the disconnected case is exactly Set Cover on a restricted but realizable class. This is a model result for a single lost supplier and one spare, not an I1 or R2 acceptance claim.

\paragraph{The operating rule.}
The tree's vertices are unit-capacity physical suppliers; each edge is a persistent duty that either endpoint can serve. Initially all duties are served. The allocator chooses one permitted vacancy $r\in S$ before the fault; no permitted vacancy is a candidate failed supplier. Scenario $i$ specifies a failed location $f_i$, a recognition time $\nu_i$, a restoration deadline $d_i\geq\nu_i$, and the first time $g_i$ at which the failed installation is proved isolated and its old effects cleared. The selected scenario's trusted information, including its calendar of authorized handoffs, is available to the allocator at recognition. It does not move the vacancy before then. Scenarios can share a failed location but differ in which permits or support windows are available. The operator allows an interruption of the failed duty until $d_i$, but no interruption of any other duty.

When the vacancy is at $u$, moving it across $uv$ to $v$ means transferring the duty on $uv$ from $v$ to $u$. For each ordered pair $(u,v)$ the calendar lists allowed starts $s$ with durations $\ell>0$, both integers. It represents actual readiness and authority, not a message that merely promises them; the final move off the failed supplier cannot start before $g_i$. The ordinary incumbent remains load-bearing until the transfer settles, the destination is reserved until then, and the allocator cannot reuse either affected capacity or abort mid-handoff. Waiting at a settled healthy vacancy is allowed without any state cost other than time. No other faults or obligations are added during repair. In particular a timer cannot infer that an old physical effect has stopped. The model takes the calendar and isolation evidence as inputs; it does not implement discovery, authentication or a distributed agreement protocol.

\paragraph{Actual assignments.}
In a tree with one vacancy $r$, every edge must be served by its endpoint farther from $r$. Indeed the first edges incident to $r$ cannot use $r$; once those duties occupy their outer endpoints the same argument continues away from $r$. This assignment is feasible and unique. After $f$ fails, the assignment ledger still retains its old duty until safe replacement. At settlement every healthy supplier is occupied and $f$ has no assigned duty, so the orientation is the one rooted at $f$. The assignments that actually change are exactly the edges of the path $r$ to $f$; other edges have both roots on the same side of the cut they induce. A single loss can therefore force $\lceil m/2\rceil$ physical transfers in the worst case, even with the best initial placement on a path of $m\geq 2$ edges whose two endpoints may fail. For this count use sufficiently generous calendars and deadlines. This is not a lower bound on CPU time or on the size of the discovery delta.

\paragraph{The exact recovery test.}
For scenario $i$, root the tree at $f_i$. Set $q_i(f_i)=d_i$; for each other vertex $u$ with parent $v$, let $q_i(u)$ be the maximum start $s$ of a calendar slot $(s,\ell)$ on $(u,v)$ for which $s+\ell\leq q_i(v)$, also requiring $s\geq g_i$ if $v=f_i$. The maximum over no slot is $-\infty$. These values give the latest time from which it is possible to restore the failed duty, from a settled healthy vacancy. To see sufficiency, wait for the selected slot and continue inductively. For necessity, any legal sequence of vacancy moves from $u$ to $f_i$ must pass along the unique path in order, even if it visits side branches or backtracks. A detour returning to any point can be replaced by waiting at that point: it restores the same binding, and the slots and old-effect clearance do not depend on the detour. Thus no other motion improves feasibility. The recurrence scans each relevant slot once per scenario. It permits later departures to finish earlier: allowing waiting accounts for that without requiring FIFO on the listed starts.

Write $B_i=\{u:q_i(u)\geq\nu_i\}$. Including $f_i$ in $B_i$ denotes the completed orientation, not a legal place to keep the spare. Every feasible $u$ has a feasible path suffix from each parent to $f_i$, starting at the common recognition time and waiting as needed. Hence $B_i$ is a connected subtree. The exact robust placement test is $S\cap\bigcap_i B_i\ne\varnothing$. The forward implication depends on making one placement before recognition; the reverse implication uses the admitted ability to select a different valid plan after recognition. Separate per-scenario schedules with different initial placements do not establish robust feasibility. Conversely, the recurrence does not apply if reaching a supplier changes which permission will be issued or if detouring changes a later obligation.

\paragraph{Short refusal certificates.}
Root the underlying tree arbitrarily and let the top of a nonempty subtree be its vertex of least depth. It is an ancestor of every vertex of that subtree. Among the tops of the $B_i$, take one of greatest depth, $p$ in $B_j$. If $B_i$ intersects $B_j$ at $z$, both tops lie on the root-to-$z$ path. Thus the top of $B_i$ is an ancestor of $p$, and path convexity puts $p$ in $B_i$. If any subtree omits $p$, report the disjoint pair $(B_j,B_i)$. Otherwise $p$ belongs to the global intersection. For a connected component $C$ of $S$, if every $B_i$ meets $C$ and the $B_i$ meet pairwise, the same argument applied to the augmented family including $C$ supplies a common placement in $C$. Thus, if no placement is possible and no two $B_i$ are disjoint, some $B_i$ entirely misses each component of $S$. Record one such scenario per component. If $S$ is empty there is nothing to check. This gives the bound $\max(2,c)$, with a single scenario enough whenever it excludes all of $S$. The case with connected $S$ can genuinely need two: stage at either of two adjacent healthy vertices on a four-vertex path, allow only an immediate one-edge transfer before a common deadline, and lose either outer endpoint. With $n$ vertices and $k$ scenarios the backward pass and component/top checks cost $O(kn+M)$ for $M$ listed slots, including the staged-membership checks. This constructs a bounded certificate, not necessarily a minimum one.

\paragraph{Sharp disconnected case and minimum explanations.}
Allow staging only at the $c$ leaves of a star. The hub fails in every scenario. In scenario $i$ forbid the sole slot at leaf $i$ and allow it at every other leaf. For an allowed leaf, dispatch at tick 1 after isolation and finish at tick 2. Each scenario rules out just one possible staging site. No subfamily of fewer than $c$ scenarios is decisive, even though the $B_i$ all meet at the hub. More generally each scenario may forbid exactly the leaves in any given subset of a Set Cover instance. In the no-disjoint-pair regime, a subfamily excludes all staging sites if and only if the components each have at least one entirely missing subtree in that subfamily, by the argument above. The star realizes arbitrary such exclusion sets. Hence the minimum-cardinality decisive subfamily problem is NP-complete in its explicit calendar-input decision form, already for this restricted star class. The reduction concerns finding a smallest explanation, not deciding whether a robust placement exists or executing the recovery.

\paragraph{Checks, comparison and boundary.}
A separate tick transition search enumerates actual assignments and in-flight destination reservations, allowing every legal transfer, not just the path. It agrees with the forward and backward procedures on 5,891 starting states from 840 generated fixtures, including calendars whose listed starts have inverted finish order. Exhaustive star controls up to eight components and 16,384 set-cover subfamily checks reproduce the two regimes. A three-transfer control fails with a one-tick delay to actual old-effect clearance despite identical travel and deadlines. The Lean file checks the deepest-top argument, the component-cover inference, and the finite-index sharpness logic; the physical model reduction is the mathematical proof above, not a Lean theorem about a plant.

The strongest direct comparator here is a temporal tree planner that runs the same backward traversal and applies the ordinary subtree Helly argument. It ties; no CPU or operational advantage over that comparator is claimed. The narrower result is about the law the planner must carry: independent alternatives are insufficient once they spend one common vacancy, while some global refusals remain short. In contrast, the project's \href{https://github.com/eternal-november/eos-research-programmes/pull/381}{timed launch-window result} intersects subsets of start times, which can have holes for one history and need many histories even in a single window. Here the convex geometry is that of \emph{locations} under the no-endogenous-change and safe-waiting hypotheses. Standard antecedents are \href{https://people.csail.mit.edu/bdean/tdsp}{time-dependent path composition}, \href{https://www.math.u-szeged.hu/~hajnal/courses/PhD_Specialis/Gavril.pdf}{subtree intersection structure}, and \href{https://www.cs.umd.edu/~gasarch/BLOGPAPERS/Karp.pdf}{Set Cover}. No priority or independent source validation is claimed. Multiple vacancies, new duties while a transfer is pending, and uncertainty persisting after recognition need new arguments. A same-source reissue with an old effect still live must keep that effect until an actual clearance; recognizing the supplier name alone is insufficient.
\end{document}
